%% Modelo mestre de currículo - layout dos PDFs de referência
%%
%% O agente deve adaptar o conteúdo à vaga sem alterar o preâmbulo,
%% as macros, o tamanho da página, as margens ou a ordem visual das seções.
%% Compilar documentos gerados com: lualatex main_<empresa>_<cargo>.tex

\documentclass[10pt,letterpaper]{article}
\usepackage[letterpaper,left=0.90in,right=0.45in,top=0.58in,bottom=0.60in]{geometry}
\usepackage[T1]{fontenc}
\usepackage[utf8]{inputenc}
\usepackage{mathptmx}
\usepackage[scaled=0.95]{helvet}
\usepackage{microtype}
\usepackage{tabularx}
\usepackage{enumitem}
\usepackage{xcolor}
\usepackage{hyperref}

\definecolor{rulegray}{RGB}{145,145,145}
\definecolor{contextgray}{RGB}{105,105,105}
\hypersetup{colorlinks=true,urlcolor=blue,linkcolor=blue,pdftitle={Candidate - CV},pdfpagemode=UseNone}
\pagestyle{empty}
\setlength{\parindent}{0pt}
\setlength{\parskip}{0pt}
\setlength{\emergencystretch}{2em}

% O cabeçalho/títulos usam Helvetica e o texto usa Times, como nos PDFs de referência.
\newcommand{\cvname}[1]{{\rmfamily\bfseries\fontsize{25}{28}\selectfont #1}\par\vspace{2pt}}
\newcommand{\contactline}[1]{{\sffamily\fontsize{9.5}{12}\selectfont #1}\par\vspace{12pt}}
\newcommand{\sectiontitle}[1]{\vspace{4pt}{\sffamily\bfseries\fontsize{15}{18}\selectfont #1}\par\vspace{1pt}{\color{rulegray}\rule{\linewidth}{0.45pt}}\par\vspace{5pt}}
\newcommand{\skillline}[2]{{\rmfamily\fontsize{10}{12}\selectfont\textbf{#1}: #2\par}}
\newcommand{\experienceentry}[3]{\begin{tabularx}{\linewidth}{@{}p{0.47\linewidth}@{}p{0.40\linewidth}@{}>{\raggedleft\arraybackslash}X@{}}{\rmfamily\bfseries #1} & {\rmfamily\itshape #2} & {\rmfamily #3}\\\end{tabularx}}
\newcommand{\educationentry}[3]{\begin{tabularx}{\linewidth}{@{}>{\bfseries}p{0.48\linewidth}@{}>{\raggedleft\arraybackslash}X@{}}#1 & #2\\\end{tabularx}\par{\rmfamily #3}\par}
\newenvironment{cvbullets}{\begin{itemize}[leftmargin=15pt,itemsep=1pt,topsep=0pt,parsep=0pt,partopsep=0pt]\rmfamily\fontsize{10}{12}\selectfont}{\end{itemize}}
\newcommand{\contextline}[1]{{\rmfamily\itshape\color{contextgray}\fontsize{9.5}{11.5}\selectfont Contexto: #1\par}}
\newcommand{\trainingline}[2]{{\rmfamily\fontsize{10}{12}\selectfont\textbf{#1}: #2\par}}

\begin{document}
\cvname{Candidate Name}
\contactline{\href{https://www.linkedin.com/in/PLACEHOLDER}{LinkedIn}\enspace | \enspace\href{https://github.com/PLACEHOLDER}{GitHub}\enspace | \enspace E-mail: YOUR.EMAIL@example.com\enspace | \enspace Telefone: +55 00 00000-0000}

\sectiontitle{Habilidades}
\skillline{Linguagens}{Java 17/21, JavaScript, TypeScript}
\skillline{Backend e APIs}{Spring Boot (projetos pessoais), REST APIs, Spring Data JPA, Flyway}
\skillline{Frontend}{React, Angular (projetos pessoais), HTML/CSS}
\skillline{Dados}{SQL, MySQL, PostgreSQL, modelagem de dados}
\skillline{Testes, DevOps e colaboração}{JUnit, Mockito, Testcontainers, Git, Maven, Gradle, Docker, Grafana, code review}
\skillline{Corporativo e integrações}{SAP (RFC, BAPI, JCo), ABAP, Mendix, OutSystems, SMTP}
\skillline{Cloud}{AWS Fundamentals em andamento via AWS Skill Builder; ECS, Auto Scaling e serviços essenciais}

\sectiontitle{Resumo}
{\rmfamily\fontsize{10}{12}\selectfont Desenvolvedor de Software com mais de 2 anos de experiência em aplicações corporativas, com foco em backend Java, integrações e entregas full stack. Atuação ponta a ponta, do levantamento com áreas de negócio e modelagem de dados à implementação de APIs, interfaces e integrações SAP. Experiência com Java 17/21, REST, SQL e Git; Spring Boot, testes automatizados e Docker aplicados em projetos pessoais. Atualmente, busco uma nova oportunidade enquanto estudo AWS e desenvolvo projetos próprios.\par}

\sectiontitle{Experiências}
\experienceentry{Alvoar Lácteos}{Junior Software Engineer}{Mai 2025 - Mar 2026}
\begin{cvbullets}
\item Conduzi uma aplicação corporativa crítica de onboarding ponta a ponta, participando das conversas com negócio, levantamento de requisitos, design system, modelagem de banco, frontend, backend e integrações; a solução reduziu etapas manuais de 8-10 para 3-5 e o onboarding de dias para horas para mais de 2.000 usuários mensais.
\item Desenvolvi integrações e componentes em Java 17/21 com REST APIs, SAP RFC/BAPI/JCo e bancos SQL (MySQL/PostgreSQL), incluindo um componente reutilizável de desserialização JSON que reduziu o esforço de mapeamento em aproximadamente 50\%.
\item Defini estruturas de dados, APIs e módulos compartilhados, trabalhando com Git, code review, Scrum/Kanban e validações manuais de fluxos; também atuei pontualmente com Grafana para consolidação de dados.
\end{cvbullets}
\contextline{Atuei no desenvolvimento de uma aplicação corporativa crítica para uma empresa de grande porte em processo de transformação digital do seu cadastro de produtores. O trabalho foi realizado dentro de um time multidisciplinar responsável por arquitetura, integrações e entrega ponta a ponta do produto.}

\vspace{5pt}
\experienceentry{Alvoar Lácteos}{Software Developer}{Dez 2023 - Mai 2025}
\begin{cvbullets}
\item Desenvolvi aplicações internas em OutSystems e componentes customizados com React/JavaScript, modernizando interfaces legadas e contribuindo para redução de 30\% no tempo de conclusão de tarefas.
\item Implementei formulários dinâmicos para substituir alterações recorrentes em formulários estáticos, reduzindo em 25\% os pedidos de modificação; trabalhei em contato direto com usuários para refinar requisitos e validar entregas.
\end{cvbullets}
\contextline{Atuei em uma empresa focada em digitalização corporativa, participando da modernização de aplicações internas e de iniciativas de transformação digital em colaboração próxima com os times de engenharia de software.}

\sectiontitle{Projeto em destaque}
{\rmfamily\fontsize{10}{12}\selectfont\textbf{DocFlow | Java 21, Spring Boot, PostgreSQL, MinIO, Flyway, Testcontainers, Docker | }\href{https://github.com/Marcos-Vinicius-F-Santos}{GitHub}\par}
Backend para cadastro, registro, armazenamento e compartilhamento de documentos, com metadados e estado persistidos em PostgreSQL e conteúdo binário em MinIO por meio de uma abstração de armazenamento independente de provedor.
\begin{cvbullets}
\item Estruturei domínio, serviços e endpoints REST com Spring Web MVC, Data JPA e Validation; usei Flyway para migrações, JUnit/Testcontainers para testes de integração e Docker Compose para infraestrutura local.
\item Documentei decisões arquiteturais com ADRs e specs, incluindo o tratamento explícito de consistência entre PostgreSQL e object storage sem assumir transação distribuída.
\item Portfólio pessoal em desenvolvimento com Java e Angular, além de projetos de estudo em React.
\end{cvbullets}

\newpage
\sectiontitle{Formação}
\educationentry{Unopar}{Jan 2022 - Jun 2023}{Tecnólogo em Design Gráfico - concluído}
\educationentry{Unopar Anhanguera}{Jul 2023 - Mai 2024}{Business Intelligence, Big Data e Analytics - Ciência de Dados - Pós-graduação Lato Sensu}
\educationentry{Unopar Anhanguera}{Dez 2023 - Jun 2024}{Estatística Aplicada - Pós-graduação Lato Sensu}

\sectiontitle{Formação complementar}
\trainingline{Java}{Criando a sua primeira aplicação - Alura (mar/2025)}
\trainingline{Microsserviços}{Padrões de projeto - Alura (ago/2025)}
\trainingline{Git e GitHub}{Compartilhando e colaborando em projetos - Alura (nov/2024)}
\trainingline{Engenharia de requisitos}{Como levantar, documentar e validar - Alura (mar/2023)}
\trainingline{UML e BPMN}{Modelagem de soluções; mapeamento de processos com Bizagi (abr/2023)}

\sectiontitle{Idiomas}
\trainingline{Inglês}{Fluente / proficiência profissional}
\trainingline{Francês}{Básico (A1)}
\trainingline{Português}{Nativo}
\end{document}
